%\ifodd\value{page}\hbox{}\newpage\fi
\chapter*{\textit{Vāk Āmbhṛṇī}: una voce femminile nel \textit{Ṛgveda}}
\addcontentsline{toc}{chapter}{\textit{Vāk Āmbhṛṇī}: una voce femminile nel \textit{Ṛgveda}}

\textit{Devī Sūkta} occupa una posizione singolare all’interno del canone vedico. Non si presenta come un’invocazione rivolta a una divinità, né come un racconto mitologico: è un inno in prima persona, in cui la Parola stessa prende la parola. La reiterazione dell’\textit{aham} non costruisce un io individuale, ma un centro di autorità cosmica che si dichiara come fondamento del mondo\endnote{1}.

Secondo la tradizione di attribuzione conservata negli \textit{Anukramaṇī}, l’inno è collegato a \textit{Vāk Āmbhṛṇī}, una \textit{ṛṣikā}, una veggente donna. Questo dato va inteso correttamente. Nel mondo vedico, l’attribuzione non coincide con l’idea moderna di autorialità biografica: il \textit{ṛṣi} non “inventa” il testo, ma lo vede. L’inno è pronunciato dalla dea \textit{Vāc}, ma è stato colto, trasmesso e reso udibile attraverso la visione di una donna\endnote{2}. La presenza di una voce femminile non è dunque un ornamento simbolico, ma una collocazione precisa dell’autorità della visione.

Il fatto che una \textit{ṛṣikā} sia associata a uno dei testi più densi e sovrani del \textit{Ṛgveda} non è secondario. L’io che parla nel \textit{Devī Sūkta} non supplica, non giustifica, non argomenta: afferma. Sostiene gli dèi, attraversa i mondi, rende possibile la conoscenza, la parola efficace e la vita stessa. In questo senso, il femminile non è tematizzato come categoria sociale, ma come forma di autorità ontologica. La donna non è qui oggetto del sacro: ne è il luogo di manifestazione.

Talvolta si è tentati di definire \textit{Vāk Āmbhṛṇī} come una \textit{“yoginī”}. Questa definizione va trattata con cautela. Il termine appartiene a una costellazione concettuale molto più tarda, legata allo sviluppo delle tradizioni yogiche e tantriche. Non esiste nel \textit{Ṛgveda} un lessico tecnico che consenta di parlare di \textit{yoga} o di pratiche iniziatiche nel senso successivo. Tuttavia, se per \textit{“yoginī”} si intende una donna che ha una realizzazione diretta e integrata della potenza della parola, allora il testo autorizza una lettura di questo tipo sul piano dell’esperienza, non della dottrina\endnote{3}.

Il \textit{Devī Sūkta} non descrive tecniche, né prescrive esercizi interiori. Mostra piuttosto una coscienza che si riconosce come principio. La Parola non è uno strumento dell’uomo per nominare il mondo: è la forza che rende possibile ogni atto, ogni conoscenza, ogni funzione vitale. In questo senso, il testo anticipa — senza formularle — alcune intuizioni che diventeranno centrali nelle tradizioni \textit{śākta} e tantriche: l’unità di parola, potere e presenza.

È importante non leggere retrospettivamente il \textit{Devī Sūkta }come un testo “tantrico \textit{ante litteram}”. La sua forza sta proprio nel non essere sistematico. Il \textit{Ṛgveda} non costruisce una teologia della Dea nel senso posteriore, ma lascia emergere una figura in cui la parola si identifica con il fondamento del reale. Proprio questa apertura, non ancora irrigidita in dottrina, renderà l’inno un punto di riferimento privilegiato per le elaborazioni successive.

La voce che parla nel \textit{Devī Sūkta} non chiede di essere creduta per autorità esterna. Dice piuttosto: “ascolta”. Molti, afferma la Dea, dimorano vicino a lei senza riconoscerla. La distanza non è spaziale, ma percettiva. In questo scarto tra prossimità e riconoscimento si colloca il gesto della veggente: non spiegare il sacro, ma rendere possibile l’ascolto. \textit{Vāk Āmbhṛṇī} non si impone come interprete: lascia che la parola parli.

In questa prospettiva, la figura della \textit{ṛṣikā} non va idealizzata né mitizzata. La sua grandezza sta nel gesto di trasmissione: prestare voce a una parola che non appartiene a nessuno e che tuttavia attraversa tutto. Il \textit{Devī Sūkta} resta così una soglia: non ancora una teologia della \textit{Devī}, ma già una dichiarazione di sovranità del femminile come principio parlante del cosmo.

\section*{Note}
%\addcontentsline{toc}{section}{Note}

\begin{enumerate}
\item J.\,A. Jamison – J.\,P. Brereton, \textit{The Rigveda: The Earliest Religious Poetry of India}, Oxford University Press, 2014, ad RV 10.125: l’inno è descritto come una \textit{ātmastuti}, una auto-lode in prima persona della dea Vāc.
\item La tradizione degli \textit{Anukramaṇī} attribuisce l’inno a \textit{Vāk Āmbhṛṇī}; cfr. Jamison–Brereton, sezione “Deities and Poets of the \textit{Rigveda”}, e intestazione a RV 10.125.
\item Sul rischio di anacronismo nell’applicazione di categorie yogiche o tantriche al Ṛgveda, si veda la cautela metodologica adottata negli studi storici sullo sviluppo delle tradizioni \textit{śākta} e yogiche; il testo vedico consente letture esperienziali, ma non dottrinali in senso tecnico.
\end{enumerate}
